\chapter{Algebraic Identities}
\index{Algebraic Identities}

% ============================================================
% FORMULA SHEET
% ============================================================
\vspace*{-1cm}
\section{Mostly Used Algebraic Identities}

\begin{formulasheet}
\begin{enumerate}
	\item $(a+b)^2= a^2+2ab+b^2 = (a-b)^2+4ab$
	\item $(a-b)^2= a^2-2ab+b^2 = (a+b)^2-4ab$
	\item $a^2+b^2= \begin{dcases}
		(a+b)^2-2ab & \\
		(a-b)^2+2ab & \\
		\dfrac{1}{2}\left\{ (a+b)^2+(a-b)^2\right\} & \\
	\end{dcases}$
	\item $\begin{dcases}
		a^2-b^2 &=  (a+b)\cdot(a-b) \\
		(a+b)^2-(a-b)^2 &= 4ab \\
		(a+b)(a-b) &= a^2-b^2 \\
		ab &= \left(\dfrac{a+b}{2}\right)^2- \left(\dfrac{a-b}{2}\right)^2\\
	\end{dcases}$
	\item $(a+b+c)^2=a^2+b^2+c^2+2ab+2bc+2ca$
	\item $a^2+b^2+c^2-ab-bc-ca= \dfrac{1}{2} \left\{(a-b)^2+(b-c)^2+(c-a)^2\right\}$
	\item $(a+b)^3=a^3+3a^2b+3ab^2+b^3 = a^3+b^3+3ab(a+b)$
	\item $(a+b)^3=a^3-3a^2b+3ab^2-b^3 = a^3-b^3-3ab(a-b)$
	\item $a^3+b^3= \begin{dcases}
		(a+b)^3-3ab(a+b) & \\
		(a+b)\left\{(a+b)^2-3ab\right\}& \\
		(a+b)(a^2-ab+b^2) &\\
	\end{dcases}$
	\item $a^3-b^3= \begin{dcases}
		(a-b)^3+3ab(a-b) & \\
		(a-b)\left\{(a-b)^2+3ab\right\}& \\
		(a-b)(a^2+ab+b^2) &\\
	\end{dcases}$
	\item $a^3+b^3+c^3-3abc = \begin{dcases}
		(a+b+c)(a^2+b^2+c^2-ab-bc-ca) &\\
		\dfrac{1}{2}(a+b+c)\left\{(a-b)^2+(b-c)^2+(c-a)^2\right\}
	\end{dcases}$
	\item $(a+b+c)^3=a^3+b^3+c^3+3(a+b)(b+c)(c+a)$
\end{enumerate}
\end{formulasheet}
\section{Some Other Algebraic Identities}
\begin{formulasheet}
\begin{enumerate}
\item $(a+b-c)^2=(c-a-b)^2=a^2+b^2+c^2+2ab-2bc-2ca$
\item $(a-b+c)^2=(b-a-c)^2=a^2+b^2+c^2-2ab-2bc+2ca$
\item $(-a+b+c)^2=(a-b-c)^2=a^2+b^2+c^2-2ab+2bc-2ca$
\item $(a+b+c)(ab+bc+ca)=(a+b)(b+c)(c+a)+abc$
\item $(a+b+c)(a+b-c)(c+a-b)(b+c-a)= 2a^2b^2+2b^2c^2+2c^2a^2-a^4-b^4-c^4$
\item \((a+b+c)^3 - (b+c-a)^3 - (a+c-b)^3 - (a+b-c)^3 = 24abc\)
\item $a^2+b^2+c^2+ab+bc+ca= \dfrac{1}{2} \left\{(a+b)^2+(b+c)^2+(c+a)^2\right\}$
\item $(a^2 + b^2)(c^2 + d^2) = (ac - bd)^2 + (ad + bc)^2$
\item $(a^2 + nb^2)(c^2 + nd^2) = (ac \pm nbd)^2 + n(ad \mp bc)^2$
\item $a^4 + 4b^4 = (a^2 + 2ab + 2b^2)(a^2 - 2ab + 2b^2)$
\item \(a^n - b^n = (a-b)(a^{n-1} + a^{n-2}b + \dots + ab^{n-2} + b^{n-1})\)
\item \(\dfrac{1}{(x-y)^{2}}+\dfrac{1}{(y-z)^{2}}+\dfrac{1}{(z-x)^{2}}=\left(\dfrac{1}{x-y}+\dfrac{1}{y-z}+\dfrac{1}{z-x}\right)^{2}\)
\item $(x + y + z)^3 - (x^3 + y^3 + z^3)=3(x + y)(y + z)(z + x)$
\item $(x + y + z)^5 - (x^5 + y^5 + z^5)=5(x + y)(y + z)(z + x)(x^2 + y^2 + z^2 + xy + yz + zx)$
\item \(\sqrt{a \pm \sqrt{b}} = \sqrt{\dfrac{a+\sqrt{a^2-b}}{2}} \pm \sqrt{\dfrac{a-\sqrt{a^2-b}}{2}}\)
\item \((a^2+b^2+c^2+d^2)(x^2+y^2+z^2+w^2) = (ax-by-cz-dw)^2 + \\ (ay+bx+cw-dz)^2 + (az-bw+cx+dy)^2 + (aw+bz-cy+dx)^2\)

\end{enumerate}
\end{formulasheet}

\section{Brahmagupta-Fibonacci Identity (Generalized)}

The generalized Brahmagupta-Fibonacci identity (also known as the Brahmagupta identity for quadratic forms) states that for any integer \(n\):

\[
(a^2 + nb^2)(c^2 + nd^2) = (ac \pm nbd)^2 + n(ad \mp bc)^2
\]

\textbf{Special cases:}
\begin{align}
\text{When } n = 1&: \quad (a^2 + b^2)(c^2 + d^2) = (ac \pm bd)^2 + (ad \mp bc)^2 \\
\text{When } n = 2&: \quad (a^2 + 2b^2)(c^2 + 2d^2) = (ac \pm 2bd)^2 + 2(ad \mp bc)^2 \\
\text{When } n = -1&: \quad (a^2 - b^2)(c^2 - d^2) = (ac \mp bd)^2 - (ad \mp bc)^2
\end{align}

\textbf{Example:} For \(a = 2\), \(b = 1\), \(c = 3\), \(d = 2\), \(n = 3\):
\[
(2^2 + 3 \cdot 1^2)(3^2 + 3 \cdot 2^2) = (2 \cdot 3 + 3 \cdot 1 \cdot 2)^2 + 3(2 \cdot 2 - 3 \cdot 1)^2
\]
\[
(4 + 3)(9 + 12) = (6 + 6)^2 + 3(4 - 3)^2
\]
\[
7 \cdot 21 = 12^2 + 3 \cdot 1^2 = 144 + 3 = 147
\]

%\newpage

\section{Euler's Four-Square Identity}

Euler's identity shows that the product of two sums of four squares is itself a sum of four squares:

\[
\begin{aligned}
&(a_1^2 + a_2^2 + a_3^2 + a_4^2)(b_1^2 + b_2^2 + b_3^2 + b_4^2) = \\
&\quad (a_1b_1 - a_2b_2 - a_3b_3 - a_4b_4)^2 + \\
&\quad (a_1b_2 + a_2b_1 + a_3b_4 - a_4b_3)^2 + \\
&\quad (a_1b_3 - a_2b_4 + a_3b_1 + a_4b_2)^2 + \\
&\quad (a_1b_4 + a_2b_3 - a_3b_2 + a_4b_1)^2
\end{aligned}
\]

Alternatively, using a more compact representation with signs:

\[
\begin{aligned}
&(a_1^2 + a_2^2 + a_3^2 + a_4^2)(b_1^2 + b_2^2 + b_3^2 + b_4^2) = \\
&\quad \sum_{i=1}^4 \left( \sum_{j=1}^4 a_j b_{p(i,j)} \right)^2
\end{aligned}
\]

where \(p(i,j)\) represents appropriate permutations with sign changes.

\textbf{Example:} Let \(a_1=1, a_2=2, a_3=0, a_4=0\) and \(b_1=1, b_2=1, b_3=1, b_4=1\):
\begin{align*}
\text{Left side} &= (1^2 + 2^2 + 0^2 + 0^2)(1^2 + 1^2 + 1^2 + 1^2) = (1+4)(4) = 20 \\
\text{Right side} &= (1\cdot1 - 2\cdot1 - 0\cdot1 - 0\cdot1)^2 + (1\cdot1 + 2\cdot1 + 0\cdot1 - 0\cdot1)^2 + \\
&\quad (1\cdot1 - 2\cdot1 + 0\cdot1 + 0\cdot1)^2 + (1\cdot1 + 2\cdot1 - 0\cdot1 + 0\cdot1)^2 \\
&= (-1)^2 + (3)^2 + (-1)^2 + (3)^2 = 1 + 9 + 1 + 9 = 20
\end{align*}

%\newpage

\section{The Difference of Fourth Powers Factorization}

The factorization of the difference of fourth powers can be expressed in several forms:

\textbf{Standard form:}
\[
a^4 - b^4 = (a - b)(a + b)(a^2 + b^2)
\]

\textbf{Extended form with \(a^2b^2\) term:}
\[
a^4 + a^2b^2 + b^4 = (a^2 + ab + b^2)(a^2 - ab + b^2)
\]

\textbf{General difference:}
\[
a^4 - b^4 = (a^2 - b^2)(a^2 + b^2) = (a - b)(a + b)(a^2 + b^2)
\]

\textbf{More general form for \(a^4 + 4b^4\) (Sophie Germain):}
\[
a^4 + 4b^4 = (a^2 + 2ab + 2b^2)(a^2 - 2ab + 2b^2)
\]

\textbf{Gaussian integer factorization:}
\[
a^4 + b^4 = (a^2 + \sqrt{2}ab + b^2)(a^2 - \sqrt{2}ab + b^2)
\]

\textbf{Examples:}
\begin{align*}
x^4 - 16 &= (x-2)(x+2)(x^2+4) \\
x^4 + x^2 + 1 &= (x^2 + x + 1)(x^2 - x + 1) \\
x^4 + 4 &= (x^2 + 2x + 2)(x^2 - 2x + 2)
\end{align*}

%\newpage

\section{The Five-term Polynomial Identity}

This remarkable identity relates the fifth power of a sum to the sum of fifth powers:

\[
\begin{aligned}
&(x + y + z)^5 - (x^5 + y^5 + z^5) = \\
&\quad 5(x + y)(y + z)(z + x)(x^2 + y^2 + z^2 + xy + yz + zx)
\end{aligned}
\]

\textbf{Expanded form:}
\[
\begin{aligned}
(x + y + z)^5 &= x^5 + y^5 + z^5 + \\
&\quad 5(x + y)(y + z)(z + x)(x^2 + y^2 + z^2 + xy + yz + zx)
\end{aligned}
\]

\textbf{Alternative symmetric representation:}
\[
\sum_{\text{sym}} x^5 + 5\sum_{\text{sym}} x^4y + 10\sum_{\text{sym}} x^3y^2 + 20\sum_{\text{sym}} x^3yz + 30x^2y^2z
\]

\textbf{Related four-term identity:}
\[
(x + y + z)^3 - (x^3 + y^3 + z^3) = 3(x + y)(y + z)(z + x)
\]

\textbf{Example:} For \(x=1, y=2, z=3\):
\begin{align*}
\text{Left side} &= (1+2+3)^5 - (1^5 + 2^5 + 3^5) = 6^5 - (1 + 32 + 243) \\
&= 7776 - 276 = 7500
\end{align*}
\begin{align*}
\text{Right side} &= 5(1+2)(2+3)(3+1)(1^2+2^2+3^2+1\cdot2+2\cdot3+3\cdot1) \\
&= 5(3)(5)(4)(1+4+9+2+6+3) \\
&= 5 \cdot 60 \cdot 25 = 5 \cdot 1500 = 7500
\end{align*}

\textbf{Verification identity:}
\[
\frac{(x+y+z)^5 - (x^5+y^5+z^5)}{(x+y)(y+z)(z+x)} = 5(x^2+y^2+z^2+xy+yz+zx)
\]

